\subsection{PASTING TOGETHER OF TOPOLOGICAL SPACES}

Let \((\mathbf{X}_{\lambda})_{\lambda \in L}\) be a family of sets, and let \(\mathbf{X}\) be the set which is the sum of the \(\mathbf{X}_{\lambda}\) (Set Theory, Chapter II, § 4, no. 8, Definition 8); we shall identify each \(\mathbf{X}_{\lambda}\) with a subset of \(\mathbf{X}\) by means of the canonical injection \(j_{\lambda}: \mathbf{X}_{\lambda} \to \mathbf{X}\) .

Let R be an equivalence relation on X such that each equivalence class of R has at most one element in each \(X_{\lambda}\) ; for each pair of indices \((\lambda, \mu)\) let \(A_{\lambda\mu}\) be the subset of \(X_{\lambda}\) consisting of the elements x for which there is an element \(y \in X_{\mu}\) which belongs to the equivalence class of x. Clearly to each \(x \in A_{\lambda\mu}\) corresponds a unique \(y \in A_{\mu\lambda}\) which is congruent to x mod R; the mappings \(h_{\mu\lambda}: A_{\lambda\mu} \to A_{\mu\lambda}\) so defined satisfy the following conditions:

\begin{enumerate}
\def\labelenumi{(\roman{enumi})}
\item
  For each \(\lambda \in L\) , \(h_{\lambda \lambda}\) is the identity mapping of \(\mathbf{A}_{\lambda \lambda} = \mathbf{X}_{\lambda}\) .
\item
  For each triple of indices \((\lambda, \mu, \nu)\) of \(\mathbf{L}\) and each \(x \in \mathbf{A}_{\lambda \mu} \cap \mathbf{A}_{\lambda \nu}\) , we have \(h_{\mu \lambda}(x) \in \mathbf{A}_{\mu \nu}\) and
\end{enumerate}

\[
h _ {\nu \lambda} (x) = h _ {\nu \mu} (h _ {\mu \lambda} (x)).\tag{1}
\]

Conversely, suppose that for each pair of indices \((\lambda, \mu)\) we are given a subset \(A_{\lambda \mu}\) of \(X_{\lambda}\) and a mapping \(h_{\mu \lambda}: A_{\lambda \mu} \to A_{\mu \lambda}\) satisfying the conditions (i) and (ii) above. It follows first of all from (ii) applied to the triples \((\lambda, \mu, \lambda)\) and \((\mu, \lambda, \mu)\) that \(h_{\lambda \mu} \circ h_{\mu \lambda}\) (resp. \(h_{\mu \lambda} \circ h_{\lambda \mu}\) ) is the restriction of \(h_{\lambda \lambda}\) (resp. \(h_{\mu \mu}\) ) to \(A_{\lambda \mu}\) (resp. \(A_{\mu \lambda}\) ); hence we deduce from (i) that \(h_{\lambda \mu}\) and \(h_{\mu \lambda}\) are bijections which are inverses of each other. Now let \(R\{x, y\}\) be the relation ``there exist \(\lambda, \mu\) such that \(x \in A_{\lambda \mu}, y \in A_{\mu \lambda}\) and \(y = h_{\mu \lambda}(x)\)''. It follows from (i) and from what precedes that \(R\) is reflexive and symmetric; on the other hand, if \(x \in A_{\lambda \mu}\) ,

\[
y = h _ {\mu \lambda} (x) \in A _ {\mu \lambda} \cap A _ {\mu \nu} \quad \text { and } \quad z = h _ {\nu \mu} (y),
\]

then also \(x = h_{\lambda \mu}(y)\) and therefore, by (ii), \(x \in A_{\lambda \mu} \cap A_{\lambda \nu}\) ; the relation (1) thus shows that \(R\) is transitive, and therefore \(R\) is an equivalence relation on \(X\) . It follows also from (i) and from the definition of \(R\) that each equivalence class mod \(R\) has at most one element in each of the sets \(X_{\lambda}\) , and that \(A_{\lambda \mu}\) is the set of all \(x \in X_{\lambda}\) for which there is an element \(y \in X_{\mu}\) congruent to \(x \mod R\) . We say that the quotient set \(X / R\) is obtained by pasting together the \(X_{\lambda}\) along the \(A_{\lambda \mu}\) by means of the bijections \(h_{\mu \lambda}\) . If \(\varphi: X \to X / R\) is the canonical mapping, the restriction of \(\varphi\) to each \(X_{\lambda}\) is a bijection of \(X_{\lambda}\) onto \(\varphi(X_{\lambda})\) .

Now suppose that each \(X_{\lambda}\) is a topological space, and let \(G_{\lambda}\) be its topology. Let G be the finest topology on the set X/R for which the mappings \(\varphi \circ j_{\lambda}\) are continuous; G is the quotient by R of the topology on X which is the sum of the topologies \(G_{\lambda}\) . We say that the topological space X/R (with the topology G) is obtained by pasting together the topological spaces \(X_{\lambda}\) along the \(A_{\lambda\mu}\) by means of the bijections \(h_{\mu\lambda}\) . The open (resp. closed) subsets of X/R are thus the canonical images of the subsets B of X which are saturated with respect to R and are such that \(B \cap X_{\lambda}\) is open (resp. closed) in \(X_{\lambda}\) for each \(\lambda \in L\) .

Since the restriction of \(\varphi\) to each \(\mathbf{X}_{\lambda}\) is a bijection onto the subset \(\mathbf{X}_{\lambda}^{\prime} = \varphi (\mathbf{X}_{\lambda})\) of \(\mathbf{X} / \mathbf{R}\) , we can transport the topology \(\mathfrak{C}_{\lambda}\) to \(\mathbf{X}_{\lambda}^{\prime}\) by means of this bijection, so that \(\mathbf{X}_{\lambda}^{\prime}\) carries a topology \(\mathfrak{C}_{\lambda}^{\prime}\) ; and the topology \(\mathfrak{C}\) on \(\mathbf{X} / \mathbf{R}\) is the finest for which the canonical injections \(\mathbf{X}_{\lambda}^{\prime} \to \mathbf{X} / \mathbf{R}\) are continuous. In general, the topology induced by \(\mathfrak{C}\) on \(\mathbf{X}_{\lambda}^{\prime}\) is coarser than \(\mathfrak{C}_{\lambda}^{\prime}\) , but not identical with the latter; even if the \(h_{\mu \lambda}\) are homeomorphisms (§ 3, Exercise 15). However, it follows from no. 4, Proposition 8 that, with the preceding notation:

PROPOSITION 9. Suppose that the \(h_{\mu \lambda}\) are homeomorphisms and that each \(\mathbf{A}_{\lambda \mu}\) is open (resp. closed) in \(\mathbf{X}_{\lambda}\) ; then each \(\varphi(\mathbf{X}_{\lambda})\) is open (resp. closed) in \(\mathbf{X}/\mathbf{R}\) and the restriction of \(\varphi\) to \(\mathbf{X}_{\lambda}\) is a homeomorphism of \(\mathbf{X}_{\lambda}\) onto the subspace \(\varphi(\mathbf{X}_{\lambda})\) of \(\mathbf{X}/\mathbf{R}\) .

